\chapter{Machine Learning Methodology}
\label{ch:ml}

\section{Data Preprocessing}

The preprocessing pipeline consists of three stages:

\subsection{Cleaning}
Missing values are filled using linear interpolation (maximum gap: 4 steps = 1 hour). Physically impossible values are clipped to engineering limits (e.g., $P_{\text{PV}} \in [0, 6000]$~W, SOC $\in [0, 100]$\%). Duplicate timestamps are removed.

\subsection{Feature Engineering}

A total of 47 features are constructed from the raw measurements:

\begin{itemize}[nosep]
    \item \textbf{Time features (12):} hour, day of week, month, is\_weekend, plus cyclical encodings ($\sin$/$\cos$) for hour, month, and day of week.
    \item \textbf{Lag features (17):} historical values at $t-1$, $t-4$, $t-96$ (24h), and $t-672$ (7 days) for load, PV, grid, SOC, temperature, and irradiance.
    \item \textbf{Rolling features (12):} rolling mean and standard deviation over 1h, 4h, and 24h windows for load, PV, temperature, and irradiance.
    \item \textbf{Derived features (6):} net power, PV capacity factor, load intensity, temperature deviation from daily mean.
\end{itemize}

All rolling and lag features use \texttt{shift(1)} to exclude the current time step, preventing data leakage.

\subsection{Train/Validation/Test Split}

The dataset is split chronologically (never shuffled):
\begin{itemize}[nosep]
    \item Training: first 70\%
    \item Validation: next 15\% (for model selection)
    \item Test: final 15\% (for reporting metrics)
\end{itemize}

The first 672 rows (7 days) are dropped due to NaN values from the 7-day lag features.

\section{Models Compared}

Four models are trained and compared:

\begin{enumerate}
    \item \textbf{Naive baseline:} Predicts the mean consumption for the same (hour, day-of-week) bin from training data.
    \item \textbf{Linear Regression:} Features are standardized (zero mean, unit variance). Provides a linear baseline.
    \item \textbf{Random Forest:} 200 trees, max depth 15. No feature scaling required.
    \item \textbf{XGBoost:} 300 trees, learning rate 0.05, max depth 8, L1/L2 regularization. Primary candidate.
\end{enumerate}

Deep learning (LSTM) is deliberately excluded because the dataset size ($\sim$34,000 rows) and feature structure (tabular with engineered lags) favor gradient-boosted trees.

\section{Evaluation Metrics}

\begin{align}
    \text{MAE} &= \frac{1}{N} \sum_{i=1}^{N} |y_i - \hat{y}_i| \\
    \text{RMSE} &= \sqrt{\frac{1}{N} \sum_{i=1}^{N} (y_i - \hat{y}_i)^2} \\
    \text{MAPE} &= \frac{100}{N} \sum_{i=1}^{N} \frac{|y_i - \hat{y}_i|}{\max(|y_i|, \epsilon)} \\
    R^2 &= 1 - \frac{\sum (y_i - \hat{y}_i)^2}{\sum (y_i - \bar{y})^2}
\end{align}

where $\epsilon = 1$~W prevents division by zero in MAPE.

\section{Reproducibility}

All models use \texttt{random\_state=42}. Dependencies are version-pinned. Results are generated from synthetic data with a fixed seed and are fully reproducible.
